% GENERATED FILE. Edit canonical lessons, then run npm run generate:books.
% canonical-source-hash: fnv1a64:fc21ea752b37a9db
% canonical-lessons: SA-C294-sanghattah, SA-C294-ghatah, SA-C294-skhalati, SA-C294-ahatah, SA-C294-murchati

\chapter{A Crash on the Road}
\label{ch:sa-a-crash-on-the-road}
\hlchaptermodality{294}

\begin{chapteropening}
\textbf{By the end of this chapter:} \emph{I can say how an accident on the road happened.}

Complete the last lesson of chapter 294: I can say how an accident on the road happened.
\end{chapteropening}

\section[\texorpdfstring{\sk{संघट्टः} (saṅghaṭṭaḥ)}{संघट्टः (saṅghaṭṭaḥ)}]{\sk{संघट्टः} (saṅghaṭṭaḥ) — a collision; a crash}

\label{lesson:SA-C294-sanghattah}



\begin{quote}
Before the new one: say the Sanskrit for a prison, then the Sanskrit for a chain.
\end{quote}

\subsection*{You'll want to know: \sk{संघट्टः}}
\textbf{\sk{संघट्टः}} — \emph{saṅghaṭṭaḥ} — \textquotedblleft{}a collision; a crash\textquotedblright{}.

\begin{quote}
\sk{मार्गे} \sk{एकः} \sk{संघट्टः} \sk{घटितः।}
\end{quote}

\begin{quote}
\emph{mārge ekaḥ saṅghaṭṭaḥ ghaṭitaḥ.}
\end{quote}

\begin{quote}
A collision happened on the road.
\end{quote}

A masculine noun.

\begin{cousinweb}
\textbf{\sk{संघट्टः}} (\emph{saṅghaṭṭaḥ}) is built from \emph{sam}, \textquotedblleft{}together\textquotedblright{}, and the root √ghaṭṭ, \textquotedblleft{}to rub, touch, shake\textquotedblright{}; MW gives \emph{saṃghaṭṭa} as a rubbing or clashing together, a collision (MW).
\end{cousinweb}

\begin{grammarlens}[title={What you've built}]
One more word for telling what happened.
\end{grammarlens}

\subsection*{Your turn}
\begin{itemize}
\raggedright
  \item \emph{Say it:} \emph{saṅghaṭṭaḥ}
  \item \emph{Say it:} \emph{saṅghaṭṭaḥ}, once more
  \item \emph{Say it:} \emph{mārge ekaḥ saṅghaṭṭaḥ ghaṭitaḥ.} — A collision happened on the road.
  \item \emph{Recall:} say the Sanskrit for a prison, then the Sanskrit for a chain, then say \emph{saṅghaṭṭaḥ} again
\end{itemize}

\begin{tcolorbox}[breakable,colback=teal!4,colframe=teal!35!black,title={Before you move on}]
What does \sk{संघट्टः} mean? (\textquotedblleft{}A collision; a crash\textquotedblright{}.) Say it once more.
\end{tcolorbox}
\section[\texorpdfstring{\sk{घातः} (ghātaḥ)}{घातः (ghātaḥ)}]{\sk{घातः} (ghātaḥ) — a blow; a bruise}

\label{lesson:SA-C294-ghatah}



\begin{quote}
Before the new one: say the Sanskrit for a chain, then the Sanskrit for a collision.
\end{quote}

\subsection*{You'll want to know: \sk{घातः}}
\textbf{\sk{घातः}} — \emph{ghātaḥ} — \textquotedblleft{}a blow; a bruise\textquotedblright{}.

\begin{quote}
\sk{चालकस्य} \sk{हस्ते} \sk{घातः} \sk{अस्ति।}
\end{quote}

\begin{quote}
\emph{cālakasya haste ghātaḥ asti.}
\end{quote}

\begin{quote}
The driver has a bruise on his hand.
\end{quote}

A masculine noun.

\begin{cousinweb}
\textbf{\sk{घातः}} (\emph{ghātaḥ}) comes from the root √han, \textquotedblleft{}to strike, kill\textquotedblright{}; MW gives \emph{ghāta} as a blow or bruise (MW).
\end{cousinweb}

\begin{grammarlens}[title={What you've built}]
One more word for telling what happened.
\end{grammarlens}

\subsection*{Your turn}
\begin{itemize}
\raggedright
  \item \emph{Say it:} \emph{ghātaḥ}
  \item \emph{Say it:} \emph{ghātaḥ}, once more
  \item \emph{Say it:} \emph{cālakasya haste ghātaḥ asti.} — The driver has a bruise on his hand.
  \item \emph{Recall:} say the Sanskrit for a chain, then the Sanskrit for a collision, then say \emph{ghātaḥ} again
\end{itemize}

\begin{tcolorbox}[breakable,colback=teal!4,colframe=teal!35!black,title={Before you move on}]
What does \sk{घातः} mean? (\textquotedblleft{}A blow; a bruise\textquotedblright{}.) Say it once more.
\end{tcolorbox}
\section[\texorpdfstring{\sk{स्खलति} (skhalati)}{स्खलति (skhalati)}]{\sk{स्खलति} (skhalati) — stumbles; slips}

\label{lesson:SA-C294-skhalati}



\begin{quote}
Before the new one: say the Sanskrit for a collision, then the Sanskrit for a blow.
\end{quote}

\subsection*{You'll want to know: \sk{स्खलति}}
\textbf{\sk{स्खलति}} — \emph{skhalati} — \textquotedblleft{}stumbles; slips\textquotedblright{}.

\begin{quote}
\sk{वृद्धः} \sk{आर्द्रे} \sk{मार्गे} \sk{स्खलति।}
\end{quote}

\begin{quote}
\emph{vṛddhaḥ ārdre mārge skhalati.}
\end{quote}

\begin{quote}
The old man stumbles on the wet road.
\end{quote}

Past (he …): \sk{सः} \sk{मार्गे} \sk{स्खलितवान्}, \emph{saḥ mārge skhalitavān}.

\begin{cousinweb}
\textbf{\sk{स्खलति}} (\emph{skhalati}) comes from the root √skhal, \textquotedblleft{}to stumble, trip, totter\textquotedblright{} (MW).
\end{cousinweb}

\begin{grammarlens}[title={What you've built}]
One more word for telling what happened.
\end{grammarlens}

\subsection*{Your turn}
\begin{itemize}
\raggedright
  \item \emph{Say it:} \emph{skhalati}
  \item \emph{Say it:} \emph{skhalati}, once more
  \item \emph{Say it:} \emph{vṛddhaḥ ārdre mārge skhalati.} — The old man stumbles on the wet road.
  \item \emph{Recall:} say the Sanskrit for a collision, then the Sanskrit for a blow, then say \emph{skhalati} again
\end{itemize}

\begin{tcolorbox}[breakable,colback=teal!4,colframe=teal!35!black,title={Before you move on}]
What does \sk{स्खलति} mean? (\textquotedblleft{}Stumbles; slips\textquotedblright{}.) Say it once more.
\end{tcolorbox}
\section[\texorpdfstring{\sk{आहतः} (āhataḥ)}{आहतः (āhataḥ)}]{\sk{आहतः} (āhataḥ) — hurt; injured}

\label{lesson:SA-C294-ahatah}



\begin{quote}
Before the new one: say the Sanskrit for a blow, then the Sanskrit for stumbles.
\end{quote}

\subsection*{You'll want to know: \sk{आहतः}}
\textbf{\sk{आहतः}} — \emph{āhataḥ} — \textquotedblleft{}hurt; injured\textquotedblright{}.

\begin{quote}
\sk{सः} \sk{वृद्धः} \sk{आहतः} \sk{अस्ति।}
\end{quote}

\begin{quote}
\emph{saḥ vṛddhaḥ āhataḥ asti.}
\end{quote}

\begin{quote}
That old man is hurt.
\end{quote}

\begin{cousinweb}
\textbf{\sk{आहतः}} (\emph{āhataḥ}) is built from \emph{ā} and the root √han, \textquotedblleft{}to strike\textquotedblright{}; MW gives \emph{āhata} as struck, beaten, hurt (MW).
\end{cousinweb}

\begin{grammarlens}[title={What you've built}]
One more word for telling what happened.
\end{grammarlens}

\subsection*{Your turn}
\begin{itemize}
\raggedright
  \item \emph{Say it:} \emph{āhataḥ}
  \item \emph{Say it:} \emph{āhataḥ}, once more
  \item \emph{Say it:} \emph{saḥ vṛddhaḥ āhataḥ asti.} — That old man is hurt.
  \item \emph{Recall:} say the Sanskrit for a blow, then the Sanskrit for stumbles, then say \emph{āhataḥ} again
\end{itemize}

\begin{tcolorbox}[breakable,colback=teal!4,colframe=teal!35!black,title={Before you move on}]
What does \sk{आहतः} mean? (\textquotedblleft{}Hurt; injured\textquotedblright{}.) Say it once more.
\end{tcolorbox}
\section[\texorpdfstring{\sk{मूर्छति} (mūrchati)}{मूर्छति (mūrchati)}]{\sk{मूर्छति} (mūrchati) — faints}

\label{lesson:SA-C294-murchati}



\begin{quote}
Before the new one: say the Sanskrit for stumbles, then the Sanskrit for injured.
\end{quote}

\subsection*{You'll want to know: \sk{मूर्छति}}
\textbf{\sk{मूर्छति}} — \emph{mūrchati} — \textquotedblleft{}faints\textquotedblright{}.

\begin{quote}
\sk{आहतः} \sk{बालः} \sk{मूर्छति।}
\end{quote}

\begin{quote}
\emph{āhataḥ bālaḥ mūrchati.}
\end{quote}

\begin{quote}
The injured boy faints.
\end{quote}

Past (he …): \sk{सः} \sk{मार्गे} \sk{मूर्छितवान्}, \emph{saḥ mārge mūrchitavān}.

\begin{cousinweb}
\textbf{\sk{मूर्छति}} (\emph{mūrchati}) comes from the root √mūrch, which MW gives as to thicken or grow rigid and also as to faint or swoon; MW spells it with one \emph{c}, \emph{mūrchati} (MW).
\end{cousinweb}

\begin{grammarlens}[title={What you've built}]
One more word for telling what happened.
\end{grammarlens}

\subsection*{Your turn}
\begin{itemize}
\raggedright
  \item \emph{Say it:} \emph{mūrchati}
  \item \emph{Say it:} \emph{mūrchati}, once more
  \item \emph{Say it:} \emph{āhataḥ bālaḥ mūrchati.} — The injured boy faints.
  \item \emph{Recall:} say the Sanskrit for stumbles, then the Sanskrit for injured, then say \emph{mūrchati} again
\end{itemize}

\begin{tcolorbox}[breakable,colback=teal!4,colframe=teal!35!black,title={Before you move on}]
What does \sk{मूर्छति} mean? (\textquotedblleft{}Faints\textquotedblright{}.) Say it once more.
\end{tcolorbox}
